\begin{ouverture}[cercleplan]{9}{Repérage et configurations}{Partie 3 · Géométrie}
{Milieux · Distances · Triangles · Aires et volumes}
\objectif{Calculer distances et milieux dans un repère orthonormé.}
\objectif{Choisir entre outils repérés et théorèmes géométriques.}
\objectif{Justifier une construction à partir de ses propriétés.}
\prerequis{Pythagore, Thalès et trigonométrie du triangle rectangle.}
\end{ouverture}
\pagestyle{otMain}\essentiel\fontsize{11}{15}\selectfont
\essSection{1}{Distances et milieux}
\begin{propriete}Dans un repère orthonormé, pour $A(x_A;y_A)$ et $B(x_B;y_B)$ :
\[AB=\sqrt{(x_B-x_A)^2+(y_B-y_A)^2}.\]
Le milieu $I$ de $[AB]$ a pour coordonnées
\[I\left(\frac{x_A+x_B}{2}\,;\,\frac{y_A+y_B}{2}\right).\]
\end{propriete}
\begin{savoirfaire}[Reconnaître une figure]
Comparer des carrés de longueurs : égalité de deux côtés pour un triangle isocèle, relation de Pythagore pour un angle droit. Pour un parallélogramme, on peut aussi montrer que les diagonales ont le même milieu ou que deux côtés opposés ont des vecteurs égaux.
\end{savoirfaire}
\begin{definition}La médiatrice de $[AB]$ est la droite perpendiculaire à $(AB)$ passant par son milieu. Ses points sont exactement les points équidistants de $A$ et de $B$.\end{definition}
\begin{attention}La formule de distance suppose des axes perpendiculaires et une même unité sur les deux axes. Un dessin déformé n’autorise pas à mesurer directement une longueur.\end{attention}
\clearpage
\essSection{2}{Réinvestir les configurations du collège}
\begin{propriete}Dans un triangle rectangle, le carré de l’hypoténuse est la somme des carrés des deux autres côtés. Réciproquement, cette égalité caractérise le triangle rectangle si l’hypoténuse présumée est le plus grand côté.\end{propriete}
\begin{savoirfaire}[Thalès]Identifier deux triangles en configuration de Thalès et établir le parallélisme requis avant d’écrire des rapports de côtés correspondants. La réciproque permet de démontrer un parallélisme si les points sont dans le même ordre.\end{savoirfaire}
\begin{propriete}Pour un angle aigu $\alpha$ d’un triangle rectangle :
\[\cos\alpha=\frac{\text{adjacent}}{\text{hypoténuse}},\quad
\sin\alpha=\frac{\text{opposé}}{\text{hypoténuse}},\quad
\tan\alpha=\frac{\text{opposé}}{\text{adjacent}}.\]
\end{propriete}
\begin{propriete}Cercle de rayon $r$ : périmètre $2\pi r$, aire du disque $\pi r^2$. Aire d’un triangle : base fois hauteur divisée par 2. Volume d’un prisme ou cylindre : aire de base fois hauteur ; pavé droit : $abc$. Pour une pyramide ou un cône, diviser aire de base fois hauteur par 3. Volume d’une boule : $4\pi r^3/3$.\end{propriete}
\otIntertitre{Choisir sa preuve}
Un repère est utile pour calculer ; une propriété géométrique peut être plus courte. Dans les corrigés, repérer ce qui relève des données, ce qui est construit et ce qui est démontré.
\clearpage\rubrique{Découverte}
\begin{blocLivre}{ex-9-1}
\begin{exercice}[1]{Distance entre deux points}
On donne $A(-1;2)$ et $B(3;5)$. Calculer $AB$.
\label{fig-ex-9-1}
\figureGeo{conf-01}
\end{exercice}
\end{blocLivre}
\begin{blocLivre}{ex-9-2}
\begin{exercice}[1]{Coordonnées d’un milieu}
Déterminer le milieu de $[AB]$ pour $A(-2;5)$ et $B(6;-1)$. Vérifier par une égalité vectorielle.
\label{fig-ex-9-2}
\figureGeo{conf-02}
\end{exercice}
\end{blocLivre}
\begin{blocLivre}{ex-9-3}
\begin{exercice}[1]{Retrouver une extrémité}
$I(2;1)$ est le milieu de $[AB]$ et $A(-1;4)$. Déterminer $B$.
\label{fig-ex-9-3}
\figureGeo{conf-03}
\end{exercice}
\end{blocLivre}
\begin{blocLivre}{ex-9-4}
\begin{exercice}[1]{Un triangle rectangle}
$A(0;0)$, $B(4;0)$ et $C(4;3)$. Montrer que le triangle est rectangle et calculer son aire.
\label{fig-ex-9-4}
\figureGeo{conf-04}
\end{exercice}
\end{blocLivre}
\begin{blocLivre}{ex-9-5}
\begin{exercice}[1]{Une médiatrice facile}
$A(-2;0)$ et $B(4;0)$. Donner le milieu de $[AB]$, puis une équation de sa médiatrice. Le point $M(1;3)$ lui appartient-il ?
\label{fig-ex-9-5}
\figureGeo{conf-05}
\end{exercice}
\end{blocLivre}

\clearpage\rubrique{Entraînement}
\begin{blocLivre}{ex-9-6}
\begin{exercice}[2]{Un triangle isocèle}
On donne $A(0;0)$, $B(4;0)$ et $C(2;3)$. Démontrer que le triangle est isocèle, puis calculer son périmètre et son aire.
\label{fig-ex-9-6}
\figureGeo{conf-06}
\end{exercice}
\end{blocLivre}
\begin{blocLivre}{ex-9-7}
\begin{exercice}[2]{Thalès dans un triangle}
Dans $ABC$, $AB=9$, $AC=12$. Les points $M\in[AB]$ et $N\in[AC]$ vérifient $AM=6$ et $(MN)\parallel(BC)$. Calculer $AN$ et le rapport des aires $AMN$ et $ABC$.
\label{fig-ex-9-7}
\figureGeo{conf-07}
\end{exercice}
\end{blocLivre}
\begin{blocLivre}{ex-9-8}
\begin{exercice}[2]{Une hauteur par trigonométrie}
Un point $P$ au sol est à 20 m du pied $B$ d’une tour verticale $BS$. L’angle $\widehat{BPS}$ vaut $35^\circ$. Exprimer la hauteur, puis l’arrondir au dixième de mètre.
\label{fig-ex-9-8}
\figureGeo{conf-08}
\end{exercice}
\end{blocLivre}
\begin{blocLivre}{ex-9-9}
\begin{exercice}[2]{Une réduction de cône}
Un cône de rayon 3 cm et hauteur 8 cm est coupé parallèlement à sa base à mi-hauteur depuis son sommet. Calculer le volume du petit cône et la fraction du volume total qu’il représente. Calculer aussi le volume d’une boule de rayon 3 cm et celui d’une pyramide de base d’aire 27 cm$^2$ et de hauteur 8 cm.
\label{fig-ex-9-9}
\figureGeo{conf-09}
\end{exercice}
\end{blocLivre}
\begin{blocLivre}{ex-9-10}
\begin{exercice}[2]{Le centre d’un parallélogramme}
$A(0;0)$, $B(4;1)$, $C(5;4)$ et $D(1;3)$. Démontrer que $ABCD$ est un parallélogramme en utilisant les milieux.
\label{fig-ex-9-10}
\figureGeo{conf-10}
\end{exercice}
\end{blocLivre}

\clearpage\rubrique{Problèmes}
\begin{blocLivre}{ex-9-11}
\begin{exercice}[3]{Un rectangle incliné}
$A(0;0)$, $B(4;2)$, $C(3;4)$ et $D(-1;2)$. Montrer que $ABCD$ est un rectangle sans utiliser le produit scalaire.
\label{fig-ex-9-11}
\figureGeo{conf-11}
\end{exercice}
\end{blocLivre}
\begin{blocLivre}{ex-9-12}
\begin{exercice}[3]{Un triangle avec un paramètre}
$A(-1;0)$, $B(3;0)$ et $C(t;2)$. Déterminer $t$ pour que $CA=CB$. Calculer alors l’aire.
\label{fig-ex-9-12}
\figureGeo{conf-12}
\end{exercice}
\end{blocLivre}
\begin{blocLivre}{ex-9-13}
\begin{exercice}[3]{Un point équidistant}
$A(0;0)$, $B(4;0)$ et $C(0;6)$. Trouver le point $O$ équidistant de ces trois points, puis le rayon du cercle qui les contient.
\label{fig-ex-9-13}
\figureGeo{conf-13}
\end{exercice}
\end{blocLivre}
\begin{blocLivre}{ex-9-14}
\begin{exercice}[3]{Un lieu de points sur une horizontale}
Les points $M(t;4)$ sont à distance 5 de l’origine $O$. Déterminer tous ces points. Donner aussi le périmètre du cercle de centre $O$ et rayon 5, puis l’aire de son disque.
\label{fig-ex-9-14}
\figureGeo{conf-14}
\end{exercice}
\end{blocLivre}
\begin{blocLivre}{ex-9-15}
\begin{exercice}[3]{Un raccourci dans un pavé}
Un pavé droit mesure 6 cm sur 8 cm, pour une hauteur de 5 cm. Calculer la diagonale de sa base, puis sa diagonale intérieure. Comparer cette dernière au trajet suivant trois arêtes.
\label{fig-ex-9-15}
\figureGeo{conf-15}
\end{exercice}
\end{blocLivre}

\clearpage
